\subsection{Moments of the first flag length}
\label{archimedean-first-length}

The preceding estimates now give the original first-length recursion.
Let $n=2r+k+1$, $k\ge1$, $m=2r+2$, and
$b\ge B\ge10^6n^7$.  Fix $T>0$ and a family indexed by
$\nu\in\mathbb Z^n$ with $\|\nu\|_1\le q_0$.  Its Harish--Chandra
parameters are $\mu_\nu=\lambda_0+\nu$.  Write
$(\sigma_\nu,H_\nu,B_\nu,d_\nu)$ for its discrete series, its
subgroup-period normalized Bessel map, and its formal degree.
All representations use the same group and subgroup Haar measures,
the same quotient measure $dq$, and the same Bessel character.
Put
\[
       F_\nu(v;g)=\sqrt{d_\nu}\,B_\nu\sigma_\nu(g)v,
       \qquad L_0(g)=\|g^{-1}f_0\|.
\]
Here $f_0$ is the original first flag vector, and $T$ is the positive
parameter in the original normalized lowering estimate, including
the case $r=0$.  The flag lengths descend to $Q$ because the
subgroup preserves the flag by unitriangular transformations.

Assume that the supplied signed dual-standard Pieri decompositions
exist at every shift used below and satisfy the parameter hypotheses
of Theorems~\ref{original-tensor-projection-bound} and
\ref{original-normalized-lowering-bound}.  The Bessel maps satisfy
the hypotheses of Theorem~\ref{original-quotient-mass-bound}, using
the original reduced gap in each shifted representation.
The level of a vector has the meaning fixed in the tensor-projection
estimate.  These assumptions retain representation existence and the
original interlacing and gap reserves; no moment estimate is assumed.

\begin{theorem}\label{original-first-length-moment-bound}
Put
\[
                 K=n^2T b^4(6B)^{2m}.
\]
If $v\in H_\nu^\infty$ has level at most $j$, $a\ge0$ is an
integer, and
\begin{equation}\label{first-length-reserve}
          \|\nu\|_1+a\le q_0,\qquad j+2a\le4n^2,
\end{equation}
then
\begin{equation}\label{first-length-moment-result}
       \int_Q L_0(g)^{2a}\|F_\nu(v;g)\|^2dq(g)
                            \le K^a\|v\|^2.
\end{equation}
In particular, the weighted integral is finite.  Moreover
$K\le T b^{6m}$.
\end{theorem}

\begin{proof}
If $B_\nu=0$, the desired bound is immediate.  We may therefore
assume $B_\nu\ne0$ whenever its lowering scalars are used.
Let $V^*$ be the dual standard representation, and use a positive
orthonormal basis $e_1,\ldots,e_n$ in the same unitary coordinates
as the tensor-projection bound.  Its action is
$(\rho(g)x)(z)=x(g^{-1}z)$.  Thus the coefficients
\[
         a_d(g)=(\rho(g)e_d)(f_0)
\]
satisfy $\sum_d|a_d(g)|^2=L_0(g)^2$.  This is evaluation in the
original flag direction, including its scale; it is not a new
choice of flag length.

Write $i_s,p_s$ for the signed Pieri inclusions and projections
from $V^*\otimes H_\nu$ to $H_{\nu-e_s}$, and let
$C_s=(\operatorname{ev}_{f_0}\otimes1)i_s$.  Under the Hermitian
identification $\omega=\operatorname{flat}(f_0)$, evaluation at
$f_0$ is the same contraction used in the normalized-adjoint lemma:
the original real Hermitian matrix satisfies $\Gamma^2=1$, so
\[
 \operatorname{flat}(f_0)=\Gamma\overline{f_0},\qquad
 \operatorname{flat}(f_0)^*\Gamma x=f_0^Tx=x(f_0).
\]
In the adapted $r=0$ frame, $f_0=\sqrt T\,e_{n-1}$ has negative
Hermitian sign; thus its flat vector is $-\sqrt T\,e_{n-1}$,
whereas evaluation contributes $+\sqrt T$.  Both original factors
are retained.
The normalized-adjoint lemma gives
$B_\nu C_s=\epsilon_s\overline{\kappa_s}B_{\nu-e_s}$, with
$\epsilon_s\in\{1,-1\}$.  Apply $B_\nu$ after evaluation at
$f_0$ to the tensor decomposition of
$\rho(g)e_d\otimes\sigma_\nu(g)v$.  Intertwining gives
\begin{equation}\label{first-length-tensor-transfer}
 a_d(g)F_\nu(v;g)
    =\sum_{s=1}^n c_s F_{\nu-e_s}(y_{ds};g),\qquad
 y_{ds}=p_s(e_d\otimes v),\qquad
 c_s=\epsilon_s\overline{\kappa_s}
                         \sqrt{d_\nu/d_{\nu-e_s}}.
\end{equation}
The scalar estimate supplies $|c_s|^2\le T b^4$.
Consequently Cauchy--Schwarz and the sum over $d$ give
\begin{equation}\label{first-length-pointwise}
 L_0^2\|F_\nu(v)\|^2
       \le nT b^4\sum_{d,s}\|F_{\nu-e_s}(y_{ds})\|^2.
\end{equation}
Zero projected vectors and zero maps at shifted indices cause no
difficulty.

Induct on $a$.  For $a=0$,
Theorem~\ref{original-quotient-mass-bound} is exactly the claim.
For the step from $a$ to $a+1$, each $y_{ds}$ has level at most
$j+2$ and
\[
       \sum_s\|y_{ds}\|^2\le(6B)^{2m}\|v\|^2.
\]
Also $\|\nu-e_s\|_1\le\|\nu\|_1+1$, so the two reserves
in \eqref{first-length-reserve} imply precisely those needed for
the induction hypothesis at $(\nu-e_s,j+2,a)$.
Multiply \eqref{first-length-pointwise} by $L_0^{2a}$ and
integrate nonnegative functions.  The induction hypothesis and the
projection sum yield
\[
 \int_Q L_0^{2(a+1)}\|F_\nu(v)\|^2dq
 \le nT b^4K^a\sum_{d,s}\|y_{ds}\|^2
 \le n^2T b^4(6B)^{2m}K^a\|v\|^2.
\]
This proves the step without a prior integrability assumption.
Finally $b\ge36$, $n^2\le b$, and $B\le b$ give
$(6B)^{2m}\le b^{3m}$ and
$n^2b^4\le b^5$.  Since $m\ge2$, $3m+5\le6m$, proving
$K\le T b^{6m}$.
\end{proof}

\begin{corollary}\label{original-flag-product-norm-bound}
Let $0\le\ell\le r+1$, and assume the derivative-energy and
flag-ratio hypotheses of Corollary~\ref{mixed-ratio-length-product}
for $v\in S_0\subset H_\nu^\infty$.  Set
\[
 a_\ell=\begin{cases}
 0,&\ell=0,\\
 1,&\ell=1,\\
 \ell+1,&\ell\ge2.
 \end{cases}
\]
If $v$ has level at most $j$ and
$\|\nu\|_1+a_\ell\le q_0$, $j+2a_\ell\le4n^2$, then
\begin{align*}
 \int_Q\prod_{i<\ell}L_i^2\,\|F_\nu(v)\|^2dq
 &\le K^\ell b^{2\ell(\ell-1)}
                           \int_Q\|F_\nu(v)\|^2dq\\
 &\le T^\ell b^{4m^2}\int_Q\|F_\nu(v)\|^2dq.
\end{align*}
For $\ell=0$ the first line is an identity without any level or
shift-reserve condition.
\end{corollary}
\begin{proof}
The first-length theorem supplies the formerly missing moment in
Corollary~\ref{mixed-ratio-length-product}.  That corollary gives
the first displayed bound with $\|v\|^2$ on the right.
If $B_\nu\ne0$, Corollary~\ref{original-quotient-mass-isometry}
identifies this with the same unweighted quotient norm.  If
$B_\nu=0$, both sides vanish directly.  The second line follows
from $K\le T b^{6m}$ and
$6m\ell+2\ell(\ell-1)\le4m^2$ for $2\ell\le m$.
\end{proof}

For an additional exterior shift of size $\ell$, the needed
integer-shift reserves are therefore $0$, $2$, and $2\ell+1$
for $\ell=0$, $\ell=1$, and $\ell\ge2$, respectively.  This
assembles the original first-length estimate and the simultaneous
flag moments \cite{first-length-original-recursion}.  The quotient
norm in the corollary is the original one.  Its identification with
the original leading-minor Jacquet norm, and the conversion of the
integer-family reserve to the original real-parameter gap budget,
remain separate obligations.  No factor $T^\ell$ is dropped.
